\mxnewpage\section*{EV2-A -- Blanc 1}\label{EV2-A}
\addcontentsline{toc}{section}{EV2-A -- Blanc 1}
\textbf{Sujet réservé au professeur}. Chapitres 1 à 7.\par Durée : 2 h. Sans calculatrice. Note sur 20. 14–16 janvier 2027.
\subsection*{Partie 1 -- Automatismes (6 points)}
Douze questions, \num{0,5} point chacune. Une seule réponse correcte ; aucune justification demandée. Une réponse incorrecte ou absente rapporte 0, sans retrait de points.
\begin{enumerate}
\item Calculer $4+2\times3$.\par A.~$13$\quad B.~$18$\quad C.~$10$\quad D.~$9$
\item Calculer $\frac12+\frac13$.\par A.~$\frac{2}{5}$\quad B.~$\frac{1}{6}$\quad C.~$\frac{5}{3}$\quad D.~$\frac{5}{6}$
\item Calculer $\dfrac{2^4}{2^{2}}$.\par A.~$2$\quad B.~$4$\quad C.~$1$\quad D.~$9$
\item Une valeur de 50 augmente de \pc{20}. Quelle est la valeur finale ?\par A.~$40$\quad B.~$10$\quad C.~$60$\quad D.~$70$
\item Une valeur de 80 est obtenue après une baisse de \pc{20}. Quelle est la valeur initiale ?\par A.~$64$\quad B.~$100$\quad C.~$96$\quad D.~$60$
\item Quel est le coefficient multiplicateur global d'une hausse de \pc{25} suivie d'une baisse de \pc{20} ?\par A.~$\frac{9}{10}$\quad B.~$\frac{4}{5}$\quad C.~$1$\quad D.~$\frac{21}{20}$
\item La solution de l'équation $2x+4=0$ est~:\par A.~$2$\quad B.~$8$\quad C.~$-8$\quad D.~$-2$
\item La pente de la droite passant par $\pt{0;4}$ et $\pt{2;10}$ est~:\par A.~$-3$\quad B.~$3$\quad C.~$2$\quad D.~$6$
\item La moyenne de $4, 6, 8$ est~:\par A.~$7$\quad B.~$18$\quad C.~$6$\quad D.~$4$
\item Un secteur de $90^\circ$ représente quelle fraction du total ?\par A.~$\frac{1}{2}$\quad B.~$\frac{1}{3}$\quad C.~$\frac{3}{4}$\quad D.~$\frac{1}{4}$
\item Une probabilité vaut $\dfrac{2}{5}$. Quelle est la probabilité de son complément ?\par A.~$\frac{1}{5}$\quad B.~$\frac{3}{5}$\quad C.~$\frac{2}{5}$\quad D.~$\frac{5}{2}$
\item Pour $f(x)=x^2+4$, quelle est la valeur de $f(-2)$ ?\par A.~$16$\quad B.~$0$\quad C.~$6$\quad D.~$8$
\end{enumerate}
\mxnewpage\subsection*{Partie 2 -- Exercices indépendants (14 points)}
Toute démarche doit être justifiée. Les résultats exacts sont attendus sauf indication contraire.
\mxneed{12}
\subsection*{Exercice 1 -- Une probabilité inconnue (5 points)}
Deux événements $A$ et $B$ d’une même expérience aléatoire vérifient $P(A)=\dfrac25$, $P(B)=\dfrac12$ et $P_A(B)=\dfrac34$. On souhaite compléter les probabilités d’un arbre dont le premier niveau distingue $A$ de son complémentaire.
\begin{enumerate}
\item Calculer $P(A\cap B)$. \hfill(1 pt)
\item Déterminer $P_{\overline A}(B)$ puis compléter un arbre. \hfill(2 pts)
\item Déterminer $P_B(A)$ et discuter l’indépendance. \hfill(2 pts)
\end{enumerate}
\mxneed{12}
\subsection*{Exercice 2 -- Position relative et variations (5 points)}
Les fonctions $f$ et $g$ sont définies sur $\intff{-1}{4}$ par $f(x)=x^3-3x^2+2$ et $g(x)=2$. Leurs courbes sont tracées dans un même repère. Il faut distinguer les variations de $f$ de la position de sa courbe par rapport à celle de $g$.
\begin{enumerate}
\item Calculer $f\prime$ et étudier les variations. \hfill(1 pt)
\item Comparer les courbes de f et g sur le domaine. \hfill(2 pts)
\item Déterminer le minimum et le maximum de $f$ sur ce domaine. \hfill(2 pts)
\end{enumerate}
\mxneed{12}
\subsection*{Exercice 3 -- Une médiatrice (4 points)}
Dans un repère orthonormé, on donne les points $A\pt{0;2}$, $B\pt{6;4}$ et $P\pt{2;6}$. On note $I$ le milieu de $[AB]$. On cherche une équation de la médiatrice de ce segment, puis on utilise cette droite pour comparer deux distances.
\begin{enumerate}
\item Donner le milieu de AB et une normale à sa médiatrice. \hfill(1 pt)
\item Déterminer l’équation de la médiatrice. \hfill(1 pt)
\item P est-il équidistant de A et B ? Vérifier par deux méthodes. \hfill(2 pts)
\end{enumerate}
